% !TeX TXS-program:compile = txs:///lualatex | txs:///view
\documentclass{article}
\usepackage[spanish]{babel}
\spanishdatedel
\usepackage{lipsum}
\input{colores} % Apunta al archivo.
\input{dict_ecuaciones}
\input{quimica}
\input{intervalos}
\title{Tomando el control de \LaTeX}
\author{Oromion}
\date{\today}
\begin{document}
\maketitle

\textcolor{myGreen}{Un texto en color verde.}

La ecuación de Fourier en línea \Fourier[inLine]{myGreen} tiene que aparecer en línea con este texto o sino del modo display \Fourier[disp]

Modo ecuación numerada \Fourier[eq]

\cuadratica{}

\Fourier[eq]

Ahoras las ecuaciones de Maxwell

\textcolor{myBlue}{\Maxwell{}}

\adrenalina{}

Ahora vamos a dibujar intervalos

\begin{tikzpicture}
\reaLine{-1}{2}
\interval{-1}{2}{0.3}{<-}{fill}{inf}{myBlue}%fill = none
\end{tikzpicture}
\end{document}